% Part: lambda-calculus
% Chapter: introduction
% Section: minimization

\documentclass[../../../include/open-logic-section]{subfiles}

\begin{document}

\olfileid{lam}{int}{min}
\olsection{الدوال \usetoken{S}{lambda definable} مغلقة تحت التصغير}

\begin{lem}
افترض أن~$f(x,y)$ !!{lambda definable}. ولتكن~$g$ معرّفة بـ
\[
g(x) \simeq \umin{y}{f(x,y)}.
\]
عندئذ تكون~$g$ !!{lambda definable}.
\end{lem}

\begin{proof}
الفكرة التقريبية كما يأتي. إذا أعطينا~$x$، استعملنا حد النقطة الثابتة
اللامبداوي~$Y$ لتعريف دالة~$h_x(n)$ تبحث عن~$y$ ابتداءً من~$n$؛ وعندئذ
لا تكون~$g(x)$ إلا~$h_x(0)$. ويمكن التعبير عن~$h_x$ بوصفها حلًا لمعادلة
نقطة ثابتة:
\[
h_x(n) \simeq
\begin{cases}
n & \text{إذا كان $f(x,n) = 0$} \\
h_x(n+1) & \text{في غير ذلك.}
\end{cases}
\]

وفيما يأتي التفاصيل. بما أن~$f$ قابلة للتعريف بلامبدا، فهي
!!{lambda defined} بحد ما~$F$. وتذكر أن لدينا أيضًا حد لامبدا~$D$ بحيث
$D(M, N, \bar{0}) \red M$ و$D(M, N, \bar{1}) \red N$. ولنثبت~$x$ مؤقتًا؛
فلكي !!{lambda define} حد ما~$h_x$ نريد إيجاد حد~$H$ (يعتمد على~$x$) يحقق
\[
H(\num{n}) \equiv D(\num{n}, H(S(\num{n})), F(x, \num{n})).
\]
ويمكننا فعل ذلك باستعمال حد النقطة الثابتة~$Y$. أولًا، ليكن~$U$ هو الحد
\[
\lambd[h][\lambd[z][D(z,(h(Sz)),F(x,z))]],
\]
ثم ليكن~$H$ هو الحد~$YU$. ولاحظ أن المتغير الحر الوحيد في~$H$ هو~$x$.
ولنبين أن~$H$ يحقق المعادلة أعلاه.

بتعريف~$Y$ لدينا
\[
H = YU \equiv U(YU) = U(H).
\]
وبخاصة، لكل عدد طبيعي~$n$، لدينا
\begin{align*}
H(\num{n}) & \equiv U(H, \num{n}) \\
& \red D(\num{n}, H(S(\num{n})), F(x, \num{n})),
\end{align*}
كما هو مطلوب. ولاحظ أنك إذا أحللت العدد~$\num{m}$ محل~$x$ في السطر
الأخير، اختزل التعبير إلى~$\num{n}$ إذا اختزل~$F(\num{m}, \num{n})$
إلى~$\num{0}$، واختزل إلى~$H(S(\num{n}))$ إذا اختزل
$F(\num{m}, \num{n})$ إلى أي عدد آخر.

ولإتمام البرهان، ليكن~$G$ هو~$\lambd[x][H(\num 0)]$. عندئذ
$G$ !!{lambda define}s~$g$؛ أي إن~$G(\num m)$، لكل~$m$، يختزل
إلى~$\overline {g(m)}$ إذا كانت~$g(m)$ معرّفة، ولا تكون له صورة سوية في
غير ذلك.
\end{proof}

\end{document}
